% !TeX root = ../main.tex
\section{Where the Difference Comes From}
\label{sec:mechanism}

The two arms differ in a loss \emph{and} in a data-collection strategy. Which carries the
gap? Preference turns out not to be the essential difference: both methods weight the
candidates of a group and step along the weighted sum --- DPO puts $+1$ on the logged winner
and $-1$ on the loser, GRPO the standardized advantage on each of its $G$ completions ---
so rescaling every group's weights to a common size puts both on one probe and lets us read
the training signal directly. Appendix~\ref{app:mechanism} gives the construction, the probe
audit that disqualifies per-iteration estimates, and the per-iteration series; two results
matter here.

\paragraph{The reward hack is visible in the gradient, and it compounds.}
The update's affirmation push --- exact, embedding-free, over every group --- rises from
$-0.006$ to $+0.086$ in GRPO and from $0.008$ to $0.103$ in PTO, so \S\ref{sec:hacking}'s
drift has a training-side measurement and not only an outcome-side inference. But that
selection pressure is an order of magnitude smaller than the drift it produces: over all
sampled candidates, affirmation cues go $0.02 \to 0.54$ (GRPO) and question marks
$0.71 \to 0.06$. Small, same-signed pressure applied each iteration to an
already-more-effusive policy compounds --- which matters for mitigation, since penalising the
selection step would barely touch a drift that lives in the generation distribution.

\paragraph{The gap is the data, not the loss.}
The arms' pooled update directions agree only weakly (cosine $0.32$, corrected for how well
each is estimated). Applying each method's weighting rule to the \emph{other's} candidate
groups isolates why:

\begin{center}\small
\begin{tabular}{@{}lr@{}}
\toprule
Comparison & cos\\
\midrule
As trained (rule \emph{and} data differ) & 0.32\\
Same groups, rule swapped (PTO's data)   & 0.91\\
Same groups, rule swapped (GRPO's data)  & 0.99\\
Same rule, each method's own data        & 0.40 / 0.32\\
\bottomrule
\end{tabular}
\end{center}

\noindent
Swapping the rule on identical groups barely moves the update; holding the rule fixed and
swapping the groups leaves the arms as far apart as when trained. \textbf{At matched $K$ and
a shared oracle, DPO and group-relative weighting extract nearly the same direction from the
same eight completions.} ``PTO vs GRPO'' is therefore a statement about \emph{exploration}
--- which conversational states each method collects candidates in --- not about the loss
family, and the design question becomes which states to visit rather than which objective to
use. One caution on our own headline: PTO's $\tau$ filter and shrinking branch count cost it
two-thirds of its training signal over the run (groups trained $782 \to 281$, against
$94$--$98\%$ retention for GRPO), so a flattening PTO curve may partly be a data-starvation
curve (Appendix~\ref{app:mechanism}, Table~\ref{tab:yield}).
